% Options for packages loaded elsewhere
\PassOptionsToPackage{unicode}{hyperref}
\PassOptionsToPackage{hyphens}{url}
\documentclass[
  ignorenonframetext,
]{beamer}
\newif\ifbibliography
\usepackage{pgfpages}
\setbeamertemplate{caption}[numbered]
\setbeamertemplate{caption label separator}{: }
\setbeamercolor{caption name}{fg=normal text.fg}
\beamertemplatenavigationsymbolsempty
% remove section numbering
\setbeamertemplate{part page}{
  \centering
  \begin{beamercolorbox}[sep=16pt,center]{part title}
    \usebeamerfont{part title}\insertpart\par
  \end{beamercolorbox}
}
\setbeamertemplate{section page}{
  \centering
  \begin{beamercolorbox}[sep=12pt,center]{section title}
    \usebeamerfont{section title}\insertsection\par
  \end{beamercolorbox}
}
\setbeamertemplate{subsection page}{
  \centering
  \begin{beamercolorbox}[sep=8pt,center]{subsection title}
    \usebeamerfont{subsection title}\insertsubsection\par
  \end{beamercolorbox}
}
% Prevent slide breaks in the middle of a paragraph
\widowpenalties 1 10000
\raggedbottom
\AtBeginPart{
  \frame{\partpage}
}
\AtBeginSection{
  \ifbibliography
  \else
    \frame{\sectionpage}
  \fi
}
\AtBeginSubsection{
  \frame{\subsectionpage}
}
\usepackage{iftex}
\ifPDFTeX
  \usepackage[T1]{fontenc}
  \usepackage[utf8]{inputenc}
  \usepackage{textcomp} % provide euro and other symbols
\else % if luatex or xetex
  \usepackage{unicode-math} % this also loads fontspec
  \defaultfontfeatures{Scale=MatchLowercase}
  \defaultfontfeatures[\rmfamily]{Ligatures=TeX,Scale=1}
\fi
\usepackage{lmodern}
\ifPDFTeX\else
  % xetex/luatex font selection
\fi
% Use upquote if available, for straight quotes in verbatim environments
\IfFileExists{upquote.sty}{\usepackage{upquote}}{}
\IfFileExists{microtype.sty}{% use microtype if available
  \usepackage[]{microtype}
  \UseMicrotypeSet[protrusion]{basicmath} % disable protrusion for tt fonts
}{}
\makeatletter
\@ifundefined{KOMAClassName}{% if non-KOMA class
  \IfFileExists{parskip.sty}{%
    \usepackage{parskip}
  }{% else
    \setlength{\parindent}{0pt}
    \setlength{\parskip}{6pt plus 2pt minus 1pt}}
}{% if KOMA class
  \KOMAoptions{parskip=half}}
\makeatother
\usepackage{longtable,booktabs,array}
\newcounter{none} % for unnumbered tables
\usepackage{calc} % for calculating minipage widths
\usepackage{caption}
% Make caption package work with longtable
\makeatletter
\def\fnum@table{\tablename~\thetable}
\makeatother
\setlength{\emergencystretch}{3em} % prevent overfull lines
\providecommand{\tightlist}{%
  \setlength{\itemsep}{0pt}\setlength{\parskip}{0pt}}
% Auto-generated by infrastructure/rendering/_pdf_latex_helpers.py::extract_math_font_preamble.
% Minimal header passed to Pandoc via -H so Beamer slides pick up the same math font
% as the combined-PDF path. Adjust the manuscript's preamble.md to change the font globally.
\usepackage{unicode-math}
\setmathfont{latinmodern-math.otf}
\usepackage{bookmark}
\IfFileExists{xurl.sty}{\usepackage{xurl}}{} % add URL line breaks if available
\urlstyle{same}
\hypersetup{
  hidelinks,
  pdfcreator={LaTeX via pandoc}}

\author{\texorpdfstring{}{}}
\date{}

\begin{document}

\section{Appendix B: Notation Reference}\label{sec:notation}

\begin{frame}[fragile,allowframebreaks]{Appendix B: Notation Reference}
This appendix collects all notation, symbols, and technical terminology
used throughout the manuscript. Entries are grouped by domain and
ordered alphabetically within each section. The \textbf{First Use}
column indicates the section where each term is first introduced or
defined.

\textbf{Sections A--K:} A Linguistic Terms; B Category Theory; C
Enriched Categories and Magnitude; D Distributional Semantics and LLMs;
E Distributional Active Inference (DAIF); F Active Inference and
Cognitive Models; G Quantum and Topological Terms; H Logical and
Type-Theoretic Terms; I AI and Communication Protocols; J Diagrammatic
Reasoning; K Notation Conventions.

Figures that colour-code case roles (category graphs, string diagrams,
appendix panels) take those colours from the shared map
\texttt{CASE\_COLORS} in \texttt{src/visualization/styles.py} (keys are
\texttt{CaseRole} names). Captions state the mapping where it matters
for reading the diagram.
\end{frame}

\begin{frame}[allowframebreaks]{Linguistic Terms}
\protect\phantomsection\label{sec:notation-linguistic}
\begin{longtable}[]{@{}
  >{\raggedright\arraybackslash}p{(\linewidth - 6\tabcolsep) * \real{0.1477}}
  >{\centering\arraybackslash}p{(\linewidth - 6\tabcolsep) * \real{0.0682}}
  >{\raggedright\arraybackslash}p{(\linewidth - 6\tabcolsep) * \real{0.7159}}
  >{\centering\arraybackslash}p{(\linewidth - 6\tabcolsep) * \real{0.0682}}@{}}
\caption{Linguistic terminology and
symbols.}\label{tbl:not-linguistic}\tabularnewline
\toprule\noalign{}
\begin{minipage}[b]{\linewidth}\raggedright
Term
\end{minipage} & \begin{minipage}[b]{\linewidth}\centering
Symbol
\end{minipage} & \begin{minipage}[b]{\linewidth}\raggedright
Definition
\end{minipage} & \begin{minipage}[b]{\linewidth}\centering
First Use
\end{minipage} \\
\midrule\noalign{}
\endfirsthead
\toprule\noalign{}
\begin{minipage}[b]{\linewidth}\raggedright
Term
\end{minipage} & \begin{minipage}[b]{\linewidth}\centering
Symbol
\end{minipage} & \begin{minipage}[b]{\linewidth}\raggedright
Definition
\end{minipage} & \begin{minipage}[b]{\linewidth}\centering
First Use
\end{minipage} \\
\midrule\noalign{}
\endhead
Ablative & ABL & Case role encoding origin, source, or cause &
\href{02_case_systems.md\#sec:case-systems}{§2} \\
Accusative & ACC & Case role encoding patient, theme, or direct object &
\href{02_case_systems.md\#sec:case-systems}{§2} \\
Active--Stative alignment & --- & Alignment in which the sole argument S
splits by agentivity &
\href{02_case_systems.md\#sec:case-systems}{§2} \\
Agent-like argument & A & The agent-like argument of a transitive clause
& \href{02_case_systems.md\#sec:case-systems}{§2} \\
Alignment functor & \(F: \mathcal{U} \to \mathcal{L}\) &
Structure-preserving map from universal to language-specific case
category & \href{02_case_systems.md\#sec:case-systems}{§2} \\
Alignment type & --- & Systematic pattern governing how S, A, P are
grouped for case marking &
\href{02_case_systems.md\#sec:case-systems}{§2} \\
Case category & \(\mathcal{C}\) & Small category whose objects are case
roles and morphisms are grammatical relations &
\href{02_case_systems.md\#sec:case-systems}{§2} \\
Case frame & --- & The set of case roles activated by a particular verb
or predicate & \href{02_case_systems.md\#sec:case-systems}{§2} \\
Case-typed noun space & \(N_{\text{NOM}}, N_{\text{ACC}}, \ldots\) &
Case-specific vector subspace in a case-enriched DisCoCat model &
\href{04_categorical_semantics.md\#sec:categorical-semantics}{§4} \\
Categorial grammar & --- & Grammar assigning each lexical item an
algebraic type encoding combinatory potential &
\href{03_categorial_grammar.md\#sec:categorial-grammar}{§3} \\
Contraction & \(a^l \cdot a \to 1\) & Pregroup reduction eliminating an
adjoint pair &
\href{03_categorial_grammar.md\#sec:categorial-grammar}{§3} \\
Dative & DAT & Case role encoding recipient, goal, or beneficiary &
\href{02_case_systems.md\#sec:case-systems}{§2} \\
Deep case & --- & Fillmore's universal semantic primitive (e.g.,
Agentive, Objective) &
\href{02_case_systems.md\#sec:case-systems}{§2} \\
Ergative--Absolutive alignment & --- & Alignment grouping S = P \(\neq\)
A & \href{02_case_systems.md\#sec:case-systems}{§2} \\
Expansion & \(1 \to a \cdot a^l\) & Pregroup expansion introducing an
adjoint pair &
\href{03_categorial_grammar.md\#sec:categorial-grammar}{§3} \\
Fluid-S & --- & Alignment in which S marking varies by context or
volition & \href{02_case_systems.md\#sec:case-systems}{§2} \\
Fluid-S functor & \(F_\theta\) & Context-dependent alignment functor
parameterized by a volition feature \(\theta\) &
\href{02_case_systems.md\#sec:case-systems}{§2} \\
Functors (Alignment) & \(F_{\text{acc}}, F_{\text{erg}}\) & Alignment
functors \(\mathcal{U} \to \mathcal{L}_{\text{acc}}\)
resp.~\(\mathcal{U} \to \mathcal{L}_{\text{erg}}\) &
\href{02_case_systems.md\#sec:case-systems}{§2} \\
Language-specific case category &
\(\mathcal{L}_{\text{acc}}, \mathcal{L}_{\text{erg}}\) & Codomain
categories for accusative vs.~ergative alignment functors &
\href{02_case_systems.md\#sec:case-systems}{§2} \\
Genitive & GEN & Case role encoding possessor or source &
\href{02_case_systems.md\#sec:case-systems}{§2} \\
Formal semantics & --- & Montague's programme: meaning assigned
compositionally via truth-conditional functions &
\href{04_categorical_semantics.md\#sec:categorical-semantics}{§4} \\
Grammatical relation & --- & Morphism in a case category relating two
case roles & \href{02_case_systems.md\#sec:case-systems}{§2} \\
Instrumental & INS & Case role encoding instrument or means &
\href{02_case_systems.md\#sec:case-systems}{§2} \\
Kāraka & --- & Pāṇini's system of deep semantic roles in Sanskrit
grammar & \href{02_case_systems.md\#sec:case-systems}{§2} \\
Left adjoint & \(a^l\) & Left adjoint type satisfying
\(a^l \cdot a \leq 1\) &
\href{03_categorial_grammar.md\#sec:categorial-grammar}{§3} \\
Locative & LOC & Case role encoding location or context &
\href{02_case_systems.md\#sec:case-systems}{§2} \\
Markedness & --- & Asymmetry in the formal complexity of paradigmatic
oppositions & \href{02_case_systems.md\#sec:case-systems}{§2} \\
Monadic semantics & --- & Song's extension using monads to model
sublexical verb-root decomposition &
\href{03_categorial_grammar.md\#sec:categorial-grammar}{§3} \\
Nominative & NOM & Case role encoding agent, experiencer, or
intransitive subject &
\href{02_case_systems.md\#sec:case-systems}{§2} \\
Nominative--Accusative alignment & --- & Alignment grouping S = A
\(\neq\) P & \href{02_case_systems.md\#sec:case-systems}{§2} \\
Passivization & --- & Syntactic transformation promoting patient to
subject position &
\href{03_categorial_grammar.md\#sec:categorial-grammar}{§3} \\
Patient-like argument & P & The patient-like argument of a transitive
clause & \href{02_case_systems.md\#sec:case-systems}{§2} \\
Pregroup grammar & --- & Grammar where types have left and right
adjoints forming a compact closed category &
\href{03_categorial_grammar.md\#sec:categorial-grammar}{§3} \\
Proto-Agent & --- & Dowty's cluster of agentive entailments (volition,
sentience, causation) &
\href{02_case_systems.md\#sec:case-systems}{§2} \\
Proto-Patient & --- & Dowty's cluster of patient entailments (change of
state, causal affectedness) &
\href{02_case_systems.md\#sec:case-systems}{§2} \\
Right adjoint & \(a^r\) & Right adjoint type satisfying
\(a \cdot a^r \leq 1\) &
\href{03_categorial_grammar.md\#sec:categorial-grammar}{§3} \\
Sole argument & S & The sole argument of an intransitive clause &
\href{02_case_systems.md\#sec:case-systems}{§2} \\
Thematic role & --- & Semantic relation between a predicate and its
argument (Agent, Patient, Goal, etc.) &
\href{02_case_systems.md\#sec:case-systems}{§2} \\
Tripartite alignment & --- & Alignment in which S \(\neq\) A \(\neq\) P
(all three distinguished) &
\href{02_case_systems.md\#sec:case-systems}{§2} \\
Verb root decomposition & --- & Analysis of a verb's internal causative
and result-state layers via a monad &
\href{03_categorial_grammar.md\#sec:categorial-grammar}{§3} \\
Vocative & VOC & Case role encoding direct address &
\href{02_case_systems.md\#sec:case-systems}{§2} \\
\bottomrule\noalign{}
\end{longtable}
\end{frame}

\begin{frame}[allowframebreaks]{Category Theory}
\protect\phantomsection\label{sec:notation-category}
\begin{longtable}[]{@{}
  >{\raggedright\arraybackslash}p{(\linewidth - 6\tabcolsep) * \real{0.1477}}
  >{\centering\arraybackslash}p{(\linewidth - 6\tabcolsep) * \real{0.0682}}
  >{\raggedright\arraybackslash}p{(\linewidth - 6\tabcolsep) * \real{0.7159}}
  >{\centering\arraybackslash}p{(\linewidth - 6\tabcolsep) * \real{0.0682}}@{}}
\caption{Category theory terminology and
symbols.}\label{tbl:not-category}\tabularnewline
\toprule\noalign{}
\begin{minipage}[b]{\linewidth}\raggedright
Term
\end{minipage} & \begin{minipage}[b]{\linewidth}\centering
Symbol
\end{minipage} & \begin{minipage}[b]{\linewidth}\raggedright
Definition
\end{minipage} & \begin{minipage}[b]{\linewidth}\centering
First Use
\end{minipage} \\
\midrule\noalign{}
\endfirsthead
\toprule\noalign{}
\begin{minipage}[b]{\linewidth}\raggedright
Term
\end{minipage} & \begin{minipage}[b]{\linewidth}\centering
Symbol
\end{minipage} & \begin{minipage}[b]{\linewidth}\raggedright
Definition
\end{minipage} & \begin{minipage}[b]{\linewidth}\centering
First Use
\end{minipage} \\
\midrule\noalign{}
\endhead
Box & --- & Node in a string diagram representing a morphism &
\href{03_categorial_grammar.md\#sec:categorial-grammar}{§3} \\
Cap & \(\eta: 1 \to a^r \otimes a\) & Unit of a compact closure
(coevaluation map) &
\href{03_categorial_grammar.md\#sec:categorial-grammar}{§3} \\
Category & \(\mathcal{C}\) & Collection of objects and morphisms with
identity and associative composition &
\href{02_case_systems.md\#sec:case-systems}{§2} \\
Classifying topos & \(\mathcal{E}_{\mathbb{T}}\) & Canonical topos:
models of \(\mathbb{T}\) in a Grothendieck topos \(\mathcal{F}\)
correspond to geometric morphisms
\(\mathcal{F} \to \mathcal{E}_{\mathbb{T}}\) &
\href{06_topos_theory.md\#sec:topos-theory}{§6} \\
Codomain & \(\text{cod}(f)\) & Target object of a morphism \(f\) &
\href{02_case_systems.md\#sec:case-systems}{§2} \\
Commutative diagram & --- & Diagram in which all directed paths with the
same start and end yield equal composites &
\href{01_introduction.md\#sec:introduction}{§1} \\
Cobordism & --- & Manifold with boundary connecting two
lower-dimensional manifolds; domain of TQFT functors &
\href{08_quantum_active_inference.md\#sec:quantum-active-inference}{§8} \\
Compact closed category & --- & Monoidal category in which every object
has a left and right dual &
\href{03_categorial_grammar.md\#sec:categorial-grammar}{§3} \\
Composition & \(g \circ f\) & Sequential application of morphisms: first
\(f\), then \(g\) & \href{02_case_systems.md\#sec:case-systems}{§2} \\
Cup & \(\varepsilon: a \otimes a^r \to 1\) & Counit of a compact closure
(evaluation map) &
\href{03_categorial_grammar.md\#sec:categorial-grammar}{§3} \\
Diagram & \(D\) & DisCoPy representation of a morphism in a monoidal
category &
\href{03_categorial_grammar.md\#sec:categorial-grammar}{§3} \\
\(\dagger\)-compact closed category & --- & Compact closed category with
a contravariant involutive endofunctor \(\dagger\); framework for
ZX-calculus &
\href{08_quantum_active_inference.md\#sec:quantum-active-inference}{§8} \\
Domain & \(\text{dom}(f)\) & Source object of a morphism \(f\) &
\href{02_case_systems.md\#sec:case-systems}{§2} \\
Fiber bundle & --- & Projection \(\pi: E \to B\) whose fibers carry
role-filler structure; topos-theoretic LoT model &
\href{06_topos_theory.md\#sec:topos-theory}{§6} \\
Functor & \(F: \mathcal{C} \to \mathcal{D}\) & Structure-preserving map
between categories & \href{02_case_systems.md\#sec:case-systems}{§2} \\
Geometric theory & \(\mathbb{T}\) & Theory axiomatized by sequents with
finite conjunctions, arbitrary disjunctions, existential quantification
& \href{06_topos_theory.md\#sec:topos-theory}{§6} \\
Geometric theories (case) &
\(\mathbb{T}_{\text{typ}}, \mathbb{T}_{\text{log}}, \mathbb{T}_{\text{dist}}, \mathbb{T}_{\text{enr}}\)
& Typological, type-logical, distributional (DisCoCat), and enriched
case theories (\autoref{sec:topos-theory}) &
\href{06_topos_theory.md\#sec:topos-theory}{§6} \\
Bridge topos & \(\mathcal{E}_{\mathbb{T}_{\text{bridge}}}\) (schematic)
& Intermediate classifying topos linking two formalizations in the
Morita bridge programme &
\href{06_topos_theory.md\#sec:topos-theory}{§6} \\
Identity morphism & \(\text{id}_A\) & Morphism from \(A\) to itself
satisfying \(f \circ \text{id} = f = \text{id} \circ f\) &
\href{02_case_systems.md\#sec:case-systems}{§2} \\
Monad & --- & Endofunctor with unit and multiplication satisfying
associativity and unit laws &
\href{03_categorial_grammar.md\#sec:categorial-grammar}{§3} \\
Monoidal category & --- & Category with a bifunctor \(\otimes\) (tensor
product) and unit object \(I\) &
\href{03_categorial_grammar.md\#sec:categorial-grammar}{§3} \\
Morita equivalence &
\(\mathcal{E}_{\mathbb{T}_1} \simeq \mathcal{E}_{\mathbb{T}_2}\) &
Equivalence of classifying toposes enabling inter-theoretic transfer &
\href{06_topos_theory.md\#sec:topos-theory}{§6} \\
Morphism & \(f: A \to B\) & Arrow between objects in a category; encodes
a relation or transformation &
\href{02_case_systems.md\#sec:case-systems}{§2} \\
Natural transformation & \(\alpha: F \Rightarrow G\) & Family of
morphisms \(\alpha_A: F(A) \to G(A)\) commuting with all morphisms;
alignment maps in semantics use \(\nu_{\text{acc}}, \nu_{\text{erg}}\);
compact-closure cap \(\eta_{\mathrm{cc}}\) is distinct from IQN
curvature \(\eta_{\mathrm{IQN}}\) &
\href{02_case_systems.md\#sec:case-systems}{§2} \\
Object & \(A, B, \ldots\) & Entities in a category; in our framework,
case roles & \href{02_case_systems.md\#sec:case-systems}{§2} \\
Presheaf & \(\hat{\mathcal{C}} = [\mathcal{C}^{op}, \textbf{Set}]\) &
Contravariant functor from \(\mathcal{C}\) to \textbf{Set} &
\href{06_topos_theory.md\#sec:topos-theory}{§6} \\
Snake equation & \((1 \otimes \varepsilon) \circ (\eta \otimes 1) = 1\)
& Zigzag identity: fundamental axiom of compact closed categories &
\href{04_categorical_semantics.md\#sec:categorical-semantics}{§4} \\
String diagram & --- & Planar graph faithfully representing morphisms in
a monoidal category &
\href{03_categorial_grammar.md\#sec:categorial-grammar}{§3} \\
Subobject classifier & \(\Omega\) & Object in a topos playing the role
of a truth-value object &
\href{06_topos_theory.md\#sec:topos-theory}{§6} \\
Swap & \(\sigma_{A,B}: A \otimes B \to B \otimes A\) & Braiding morphism
permuting two objects in a symmetric monoidal category &
\href{03_categorial_grammar.md\#sec:categorial-grammar}{§3} \\
Tensor product & \(A \otimes B\) & Monoidal product representing
parallel composition or concatenation &
\href{03_categorial_grammar.md\#sec:categorial-grammar}{§3} \\
Topos & \(\mathcal{E}\) & Category with products, exponentials, and a
subobject classifier; generalized universe of sets &
\href{06_topos_theory.md\#sec:topos-theory}{§6} \\
Universal construction & --- & Object or morphism characterized by a
universal property (product, coproduct, limit) &
\href{06_topos_theory.md\#sec:topos-theory}{§6} \\
Wire & --- & Edge in a string diagram representing a type (object) &
\href{03_categorial_grammar.md\#sec:categorial-grammar}{§3} \\
\bottomrule\noalign{}
\end{longtable}
\end{frame}

\begin{frame}[fragile,allowframebreaks]{Enriched Categories and
Magnitude}
\protect\phantomsection\label{sec:notation-enriched}
\begin{longtable}[]{@{}
  >{\raggedright\arraybackslash}p{(\linewidth - 6\tabcolsep) * \real{0.1477}}
  >{\centering\arraybackslash}p{(\linewidth - 6\tabcolsep) * \real{0.0682}}
  >{\raggedright\arraybackslash}p{(\linewidth - 6\tabcolsep) * \real{0.7159}}
  >{\centering\arraybackslash}p{(\linewidth - 6\tabcolsep) * \real{0.0682}}@{}}
\caption{Enriched categories and magnitude
terminology.}\label{tbl:not-enriched}\tabularnewline
\toprule\noalign{}
\begin{minipage}[b]{\linewidth}\raggedright
Term
\end{minipage} & \begin{minipage}[b]{\linewidth}\centering
Symbol
\end{minipage} & \begin{minipage}[b]{\linewidth}\raggedright
Definition
\end{minipage} & \begin{minipage}[b]{\linewidth}\centering
First Use
\end{minipage} \\
\midrule\noalign{}
\endfirsthead
\toprule\noalign{}
\begin{minipage}[b]{\linewidth}\raggedright
Term
\end{minipage} & \begin{minipage}[b]{\linewidth}\centering
Symbol
\end{minipage} & \begin{minipage}[b]{\linewidth}\raggedright
Definition
\end{minipage} & \begin{minipage}[b]{\linewidth}\centering
First Use
\end{minipage} \\
\midrule\noalign{}
\endhead
Base-change functor & --- & Conceptual tower
\(\mathbf{Bool} \hookrightarrow [0,1] \hookrightarrow \mathbf{R}_{\geq 0}\)
showing progressively richer enrichments of grammatical categories (not
computationally instantiated; this paper implements only
\([0,1]\)-enrichment) &
\href{03_categorial_grammar.md\#sec:categorial-grammar}{§3} \\
Categorical magnitude & \(\lvert X \rvert\) & Sum \(\sum_i w_i\) where
\((w_1, \ldots, w_n)\) solves \(Z \mathbf{w} = \mathbf{1}\); measures
effective number of objects &
\href{05_enriched_categories.md\#sec:enriched-categories}{§5} \\
Composition inequality &
\(\mathcal{V}(A,B) \cdot \mathcal{V}(B,C) \leq \mathcal{V}(A,C)\) &
Enriched analogue of composition on \(([0,1], \cdot, 1)\): composite
hom-value is at least the product of intermediates &
\href{05_enriched_categories.md\#sec:enriched-categories}{§5} \\
Enriched category & \(\mathcal{V}\text{-Cat}\) & Category whose hom-sets
are replaced by objects of a monoidal category \(\mathcal{V}\) &
\href{05_enriched_categories.md\#sec:enriched-categories}{§5} \\
Enriched functor & --- & Structure-preserving map between enriched
categories respecting hom-values &
\href{05_enriched_categories.md\#sec:enriched-categories}{§5} \\
Hom-value & \(\mathcal{V}(A,B) \in [0,1]\) & Enriched analogue of
hom-set; measures degree of relation between objects &
\href{05_enriched_categories.md\#sec:enriched-categories}{§5} \\
Identity axiom & \(\mathcal{V}(A,A) = 1\) & Every object has maximal
self-relatedness in the enriched hom (equality, not mere inequality, in
our \([0,1]\)-convention) &
\href{05_enriched_categories.md\#sec:enriched-categories}{§5} \\
Similarity matrix & \(Z_{ij} = \mathcal{V}(A_i, A_j)\) & Matrix of all
pairwise hom-values; used to compute categorical magnitude &
\href{05_enriched_categories.md\#sec:enriched-categories}{§5} \\
Lawvere metric space & --- & Category enriched over
\(([0,\infty], +, 0)\); generalizes metric spaces via enriched
categories &
\href{05_enriched_categories.md\#sec:enriched-categories}{§5} \\
Magnitude homology & \(\text{MH}_n(\mathcal{C})\) & Graded homological
invariant categorifying magnitude; detects higher-dimensional holes in
enriched categories &
\href{05_enriched_categories.md\#sec:enriched-categories}{§5} \\
Morphism weight (precision) & \(w_f\), \(w_k\), \(w_c\) & Enriched
weights in \([0,1]\); subscript denotes morphism or role in VMP/DPE
(\autoref{sec:daif-results}). \textbf{Convention}: \(w_f\) always refers
to the \emph{enriched} hom-value \(\mathcal{C}(A,B)\) in the
\([0,1]\)-enriched category
(\href{05_enriched_categories.md\#sec:enriched-categories}{§5}), not the
unit-weight morphism scalar carried by the \texttt{CaseCategory} in
\href{02_case_systems.md\#sec:case-systems}{§2} --- the two number
systems are intentionally decoupled in the implementation (see
\href{05_enriched_categories.md\#sec:enriched-categories}{§5}
Architectural note). &
\href{07_cognitive_integration.md\#sec:cognitive-integration}{§7} \\
POVM element & \(E_c\) & Positive operator-valued measure element for
case role \(c\); \(P(c \mid \rho) = \text{Tr}(E_c \rho)\) &
\href{08_quantum_active_inference.md\#sec:quantum-active-inference}{§8} \\
Prediction error & \(\text{PE}(f)\) &
\(\propto w_f \cdot \lvert\mu_{\text{predicted}} - \mu_{\text{observed}}\rvert\);
case violation signal scaling with morphism weight &
\href{07_cognitive_integration.md\#sec:cognitive-integration}{§7} \\
Shannon entropy & \(H\) & Information-theoretic measure characterized as
the unique derivation of a topological operad (Bradley) &
\href{05_enriched_categories.md\#sec:enriched-categories}{§5} \\
Topological operad & --- & Operad with topological structure whose
derivations connect magnitude to entropy &
\href{05_enriched_categories.md\#sec:enriched-categories}{§5} \\
Weight vector & \(\mathbf{w}\) & Solution to
\(Z\mathbf{w} = \mathbf{1}\); entries are the effective weights of each
object &
\href{05_enriched_categories.md\#sec:enriched-categories}{§5} \\
\bottomrule\noalign{}
\end{longtable}
\end{frame}

\begin{frame}[allowframebreaks]{Distributional Semantics and LLMs}
\protect\phantomsection\label{sec:notation-distributional}
\begin{longtable}[]{@{}
  >{\raggedright\arraybackslash}p{(\linewidth - 6\tabcolsep) * \real{0.1477}}
  >{\centering\arraybackslash}p{(\linewidth - 6\tabcolsep) * \real{0.0682}}
  >{\raggedright\arraybackslash}p{(\linewidth - 6\tabcolsep) * \real{0.7159}}
  >{\centering\arraybackslash}p{(\linewidth - 6\tabcolsep) * \real{0.0682}}@{}}
\caption{Distributional semantics and LLM
terminology.}\label{tbl:not-distributional}\tabularnewline
\toprule\noalign{}
\begin{minipage}[b]{\linewidth}\raggedright
Term
\end{minipage} & \begin{minipage}[b]{\linewidth}\centering
Symbol
\end{minipage} & \begin{minipage}[b]{\linewidth}\raggedright
Definition
\end{minipage} & \begin{minipage}[b]{\linewidth}\centering
First Use
\end{minipage} \\
\midrule\noalign{}
\endfirsthead
\toprule\noalign{}
\begin{minipage}[b]{\linewidth}\raggedright
Term
\end{minipage} & \begin{minipage}[b]{\linewidth}\centering
Symbol
\end{minipage} & \begin{minipage}[b]{\linewidth}\raggedright
Definition
\end{minipage} & \begin{minipage}[b]{\linewidth}\centering
First Use
\end{minipage} \\
\midrule\noalign{}
\endhead
Attention mechanism & --- & Transformer component computing weighted
relevance between token positions &
\href{04_categorical_semantics.md\#sec:categorical-semantics}{§4} \\
Attention weight & \(\alpha_{ij}\) & Softmax-normalized score encoding
contextual relevance of token \(j\) to token \(i\) &
\href{04_categorical_semantics.md\#sec:categorical-semantics}{§4} \\
BERT & --- & Bidirectional Encoder Representations from Transformers;
masked language model producing contextualized embeddings &
\href{04_categorical_semantics.md\#sec:categorical-semantics}{§4} \\
Contextualized embedding & \(\mathbf{v}_w^{(c)}\) & Vector
representation of word \(w\) that varies with linguistic context \(c\) &
\href{04_categorical_semantics.md\#sec:categorical-semantics}{§4} \\
Compact closure map & \(\varepsilon_N: N \otimes N \to \mathbb{R}\) &
Inner product implementing pregroup contraction in the vector space
semantics &
\href{04_categorical_semantics.md\#sec:categorical-semantics}{§4} \\
Cosine similarity & \(\cos(\mathbf{u}, \mathbf{v})\) & Similarity
measure
\(\frac{\mathbf{u} \cdot \mathbf{v}}{\|\mathbf{u}\| \|\mathbf{v}\|}\)
between vectors &
\href{05_enriched_categories.md\#sec:enriched-categories}{§5} \\
Distributional Memory & --- & Baroni--Lenci tensor-based framework
structuring co-occurrence as (word, relation, word) triples &
\href{04_categorical_semantics.md\#sec:categorical-semantics}{§4} \\
DisCoCat & --- & Distributional Compositional Categorical model;
monoidal functor from pregroup grammars to vector spaces &
\href{03_categorial_grammar.md\#sec:categorial-grammar}{§3} \\
DisCoCirc & --- & Discourse-level extension of DisCoCat with persistent
entity wires &
\href{04_categorical_semantics.md\#sec:categorical-semantics}{§4} \\
Distributional hypothesis & --- & The thesis that words occurring in
similar contexts have similar meanings (Harris 1954, Firth 1957) &
\href{04_categorical_semantics.md\#sec:categorical-semantics}{§4} \\
GloVe & --- & Global Vectors for Word Representation; log-bilinear model
of word co-occurrence statistics &
\href{04_categorical_semantics.md\#sec:categorical-semantics}{§4} \\
GPT & --- & Generative Pre-trained Transformer; autoregressive language
model &
\href{04_categorical_semantics.md\#sec:categorical-semantics}{§4} \\
Meaning functor & \(F: \mathbf{Preg} \to \mathbf{FVect}\) & Monoidal
functor assigning vector spaces to types and linear maps to derivations
& \href{04_categorical_semantics.md\#sec:categorical-semantics}{§4} \\
Noun space & \(N\) & Vector space to which noun types \(n\) map under
the meaning functor &
\href{04_categorical_semantics.md\#sec:categorical-semantics}{§4} \\
Parameterized optic & --- & Categorical construction (Gavranović)
modeling attention heads as functorial lenses &
\href{09_ai_implications.md\#sec:ai-implications}{§9} \\
Self-attention &
\(\text{Attn}(Q,K,V) = \mathrm{softmax}(\frac{QK^\top}{\sqrt{d_k}})V\) &
Core transformer operation computing contextualized representations &
\href{04_categorical_semantics.md\#sec:categorical-semantics}{§4} \\
Sentence space & \(S\) & Vector space to which sentence type \(s\) maps
under the meaning functor &
\href{04_categorical_semantics.md\#sec:categorical-semantics}{§4} \\
Sentence vector & \(\overrightarrow{\text{sentence}}\) & Vector in \(S\)
computed by tensoring and contracting word meanings via DisCoCat &
\href{04_categorical_semantics.md\#sec:categorical-semantics}{§4} \\
Static embedding & \(\mathbf{v}_w\) & Fixed vector representation of
word \(w\) independent of context (e.g., Word2Vec, GloVe) &
\href{04_categorical_semantics.md\#sec:categorical-semantics}{§4} \\
Transformer & --- & Neural architecture using self-attention and
feed-forward layers for sequence processing &
\href{04_categorical_semantics.md\#sec:categorical-semantics}{§4} \\
Word2Vec & --- & Neural embedding model (Mikolov et al.~2013) learning
word vectors from local context windows &
\href{04_categorical_semantics.md\#sec:categorical-semantics}{§4} \\
Word vector & \(\mathbf{v}_w \in \mathbb{R}^d\) & \(d\)-dimensional
real-valued vector encoding distributional properties of word \(w\) &
\href{04_categorical_semantics.md\#sec:categorical-semantics}{§4} \\
\bottomrule\noalign{}
\end{longtable}
\end{frame}

\begin{frame}[fragile,allowframebreaks]{Distributional Active Inference
(DAIF)}
\protect\phantomsection\label{sec:notation-daif}
\begin{longtable}[]{@{}
  >{\raggedright\arraybackslash}p{(\linewidth - 6\tabcolsep) * \real{0.1477}}
  >{\centering\arraybackslash}p{(\linewidth - 6\tabcolsep) * \real{0.0682}}
  >{\raggedright\arraybackslash}p{(\linewidth - 6\tabcolsep) * \real{0.7159}}
  >{\centering\arraybackslash}p{(\linewidth - 6\tabcolsep) * \real{0.0682}}@{}}
\caption{Distributional Active Inference (DAIF)
terminology.}\label{tbl:not-daif}\tabularnewline
\toprule\noalign{}
\begin{minipage}[b]{\linewidth}\raggedright
Term
\end{minipage} & \begin{minipage}[b]{\linewidth}\centering
Symbol
\end{minipage} & \begin{minipage}[b]{\linewidth}\raggedright
Definition
\end{minipage} & \begin{minipage}[b]{\linewidth}\centering
First Use
\end{minipage} \\
\midrule\noalign{}
\endfirsthead
\toprule\noalign{}
\begin{minipage}[b]{\linewidth}\raggedright
Term
\end{minipage} & \begin{minipage}[b]{\linewidth}\centering
Symbol
\end{minipage} & \begin{minipage}[b]{\linewidth}\raggedright
Definition
\end{minipage} & \begin{minipage}[b]{\linewidth}\centering
First Use
\end{minipage} \\
\midrule\noalign{}
\endhead
Return distribution & \(Z(s,a)\) & Full probability distribution over
discounted cumulative rewards from state \(s\) under action \(a\) &
\href{07c_daif_results.md\#sec:daif-results}{§7c} \\
Distributional Bellman operator & \(\mathcal{T}^{\pi}\) & Distributional
analogue of the Bellman operator:
\(\mathcal{T}^{\pi}Z \stackrel{d}{=} R + \gamma Z'\) &
\href{07c_daif_results.md\#sec:daif-results}{§7c} \\
Push-forward measure & \(\mathbf{S}_{\#}\mathbb{P}\) & Path-level
push-forward of the trajectory measure (composed with decoder
\(f: \mathcal{S} \to \mathcal{X}\) in \autoref{eq:eq-7-1}); generic
push-forward written \(T_{\#}\mu\) &
\href{07c_daif_results.md\#sec:daif-results}{§7c} \\
Discount factor & \(\gamma \in (0,1)\) & Exponential time-discount in
the distributional Bellman return &
\href{07c_daif_results.md\#sec:daif-results}{§7c} \\
Quantile level & \(\tau \in [0,1]\) & Quantile index in QR-DQN and IQN;
sampled uniformly at inference time &
\href{07c_daif_results.md\#sec:daif-results}{§7c} \\
Quantile Huber loss & \(\rho_{\tau}^{\kappa}{(u)}\) & Loss function for
quantile regression: interpolates \(L^1\) and \(L^2\) via threshold
\(\kappa{}\) & \href{07c_daif_results.md\#sec:daif-results}{§7c} \\
Wasserstein distance & \(W_p(P,Q)\) & \(p\)-Wasserstein distance between
return distributions; \(W_1\) = area between CDFs &
\href{07c_daif_results.md\#sec:daif-results}{§7c} \\
Bethe free energy & \(F_{\text{Bethe}}[\mathbf{q}]\) & Tractable
approximation to variational free energy in belief propagation &
\href{07c_daif_results.md\#sec:daif-results}{§7c} \\
Expected information gain & \(\text{EIG}{(o)}\) &
\(D_{\mathrm{KL}}{(q{(\mathbf{c} \mid o)} \Vert p{(\mathbf{c})})}\);
epistemic value of observation \(o{}\) &
\href{07c_daif_results.md\#sec:daif-results}{§7c} \\
Expected free energy & \(G{(\pi)}\) & Pragmatic + epistemic + risk cost
of policy \(\pi{}\); minimised for action selection &
\href{07c_daif_results.md\#sec:daif-results}{§7c} \\
Distributional prediction error & \(\mathrm{DPE}\) & Precision-weighted
Wasserstein mismatch:
\(w_f \cdot W_1{(Z_{\text{pred}}, Z_{\text{obs}})}\) &
\href{07c_daif_results.md\#sec:daif-results}{§7c} \\
ERP amplitude profile & \(\mathrm{ERPProfile}\) & Dataclass holding
N400/P600 amplitudes, peak latencies, and time-series waveforms for each
case role & \href{07c_daif_results.md\#sec:daif-results}{§7c} \\
Return entropy & \(H[Z]\) & Shannon entropy of the return distribution:
\(-\sum_i p_i \log p_i\) &
\href{07c_daif_results.md\#sec:daif-results}{§7c} \\
Quantile calibration error & \(\mathrm{CE}\) & Mean absolute deviation
between nominal quantile levels and empirical coverage &
\href{07c_daif_results.md\#sec:daif-results}{§7c} \\
IQN risk distortion & \(\psi_{\mathrm{IQN}}(\tau)\) & Maps quantile
levels \(\tau{}\); piecewise formulas (neutral, power, tail, CVaR) match
\texttt{implicit\_quantile\_network\_update()} and the
\href{07c_daif_results.md\#sec:daif-results}{§7c} table; distinct from
\(\beta_{\mathrm{risk}}\) in \(G(\pi)\) &
\href{07c_daif_results.md\#sec:daif-results}{§7c} \\
IQN curvature & \(\eta_{\mathrm{IQN}}\) & Default \(0.71\) in code; not
the compact-closure cap \(\eta_{\mathrm{cc}}\) &
\href{07c_daif_results.md\#sec:daif-results}{§7c} \\
CVaR scale & \(\alpha_{\mathrm{CVaR}}\) & IQN mode uses
\(\psi_{\mathrm{IQN}}(\tau)=\tau\cdot\alpha_{\mathrm{CVaR}}\); default
\(0.25\) in code (distinct from softmax \(\alpha_{\mathrm{pol}}\)) &
\href{07c_daif_results.md\#sec:daif-results}{§7c} \\
Risk sensitivity (EFE) & \(\beta_{\mathrm{risk}}\) & Non-negative
coefficient on \(\text{risk}(\pi)\) in \(G(\pi)\);
\(\beta_{\mathrm{risk}}=0\) recovers standard EFE &
\href{07c_daif_results.md\#sec:daif-results}{§7c} \\
Inverse temperature (policy) & \(\alpha_{\mathrm{pol}}\) & Softmax
sharpness over policies in \(P(\pi)\) &
\href{07c_daif_results.md\#sec:daif-results}{§7c} \\
Belief precision (VMP) &
\(\Lambda_{\text{post}}, \Lambda_{\text{prior}}, \Lambda_{\text{lik}}\)
& Posterior, prior, and likelihood precision matrices; \(\Delta\Lambda\)
is their positive update magnitude &
\href{07c_daif_results.md\#sec:daif-results}{§7c} \\
Violation severity & \(S_{\text{violation}}\) & \(\in \{0, 0.5, 1.0\}\)
in ERP decomposition &
\href{07c_daif_results.md\#sec:daif-results}{§7c} \\
TD error (quantile) & \(\delta_{ij}\) & Temporal-difference residual in
QR-DQN (\autoref{eq:eq-7c-qr}) &
\href{07c_daif_results.md\#sec:daif-results}{§7c} \\
Quantile count (current) & \(N\) & Number of current-network quantile
levels in QR-DQN (Eq. \ref{eq:eq-7c-qr}); equal to the length of
\texttt{current\_quantiles} passed to \texttt{quantile\_td\_update()} &
\href{07c_daif_results.md\#sec:daif-results}{§7c} \\
Quantile count (target) & \(N'\) & Number of target-network quantile
samples in QR-DQN (Eq. \ref{eq:eq-7c-qr}) &
\href{07c_daif_results.md\#sec:daif-results}{§7c} \\
Huber threshold & \(\kappa\) & Cut-point between quadratic and linear
regimes of the Huber loss \(L_\kappa(u)\) (Eq. \ref{eq:eq-7c-huber});
default \(\kappa=1\) in \texttt{quantile\_td\_update()} &
\href{07c_daif_results.md\#sec:daif-results}{§7c} \\
Atom count (C51) & \(N_{\text{atoms}}\) & Number of support atoms
\(z_i\) in the C51 categorical representation (Eq. \ref{eq:eq-7c-c51});
default \(N_{\text{atoms}}=51\) in
\texttt{categorical\_return\_distribution()} &
\href{07c_daif_results.md\#sec:daif-results}{§7c} \\
Support bounds (C51) & \(V_{\min},\,V_{\max}\) & Lower and upper
endpoints of the atomic support spanned by \(\{z_i\}\) in the C51
representation & \href{07c_daif_results.md\#sec:daif-results}{§7c} \\
Bin count (entropy) & \(N_{\text{bins}}\) & Number of equal-width bins
used to discretise the quantile-parameterised return when computing
\(H[Z]\) (Eq. \ref{eq:eq-7c-entropy}); default \(N_{\text{bins}}=50\) in
\texttt{return\_distribution\_entropy()} &
\href{07c_daif_results.md\#sec:daif-results}{§7c} \\
\bottomrule\noalign{}
\end{longtable}
\end{frame}

\begin{frame}[allowframebreaks]{Active Inference and Cognitive Models}
\protect\phantomsection\label{sec:notation-cognitive}
\begin{longtable}[]{@{}
  >{\raggedright\arraybackslash}p{(\linewidth - 6\tabcolsep) * \real{0.1477}}
  >{\centering\arraybackslash}p{(\linewidth - 6\tabcolsep) * \real{0.0682}}
  >{\raggedright\arraybackslash}p{(\linewidth - 6\tabcolsep) * \real{0.7159}}
  >{\centering\arraybackslash}p{(\linewidth - 6\tabcolsep) * \real{0.0682}}@{}}
\caption{Active inference and cognitive modeling
terminology.}\label{tbl:not-cognitive}\tabularnewline
\toprule\noalign{}
\begin{minipage}[b]{\linewidth}\raggedright
Term
\end{minipage} & \begin{minipage}[b]{\linewidth}\centering
Symbol
\end{minipage} & \begin{minipage}[b]{\linewidth}\raggedright
Definition
\end{minipage} & \begin{minipage}[b]{\linewidth}\centering
First Use
\end{minipage} \\
\midrule\noalign{}
\endfirsthead
\toprule\noalign{}
\begin{minipage}[b]{\linewidth}\raggedright
Term
\end{minipage} & \begin{minipage}[b]{\linewidth}\centering
Symbol
\end{minipage} & \begin{minipage}[b]{\linewidth}\raggedright
Definition
\end{minipage} & \begin{minipage}[b]{\linewidth}\centering
First Use
\end{minipage} \\
\midrule\noalign{}
\endhead
Active inference & --- & Framework in which perception and action are
unified as variational free energy minimization &
\href{07_cognitive_integration.md\#sec:cognitive-integration}{§7} \\
Active sampling & --- & Agent's selection of actions to confirm or
update case assignments via sensory evidence &
\href{07_cognitive_integration.md\#sec:cognitive-integration}{§7} \\
Belief updating & --- & Bayesian posterior computation: revising
generative model parameters given new observations &
\href{07_cognitive_integration.md\#sec:cognitive-integration}{§7} \\
CEREBRUM & --- & Case-Enabled Reasoning Engine with Bayesian
Representations for Unified Modeling; treats AI models as case-bearing
entities &
\href{07_cognitive_integration.md\#sec:cognitive-integration}{§7} \\
Distributional Active Inference (DAIF) & --- & Extension replacing
scalar value summaries with full return distributions in active
inference &
\href{07_cognitive_integration.md\#sec:cognitive-integration}{§7} \\
Distributional RL & --- & Reinforcement learning operating on full
return distributions rather than scalar expected values &
\href{07_cognitive_integration.md\#sec:cognitive-integration}{§7} \\
Free energy & \(F\) & Variational bound on surprisal;
\(F = D_{KL}[q(\theta) \parallel p(\theta \mid o)] - \ln p(o)\) &
\href{07_cognitive_integration.md\#sec:cognitive-integration}{§7} \\
Free energy principle (FEP) & --- & The principle that self-organizing
systems maintain themselves by minimizing surprisal &
\href{07_cognitive_integration.md\#sec:cognitive-integration}{§7} \\
Garden-path reanalysis & --- & Restructuring of case assignments when
incoming evidence contradicts the current parse &
\href{07_cognitive_integration.md\#sec:cognitive-integration}{§7} \\
Generative model & \(p(o, s)\) & Joint probability model over
observations \(o\) and hidden states \(s\) &
\href{07_cognitive_integration.md\#sec:cognitive-integration}{§7} \\
Markov blanket & --- & Statistical boundary separating internal states
from external environment; defines agent boundary &
\href{07_cognitive_integration.md\#sec:cognitive-integration}{§7} \\
N400 & --- & Event-related brain potential peaking \textasciitilde400ms
post-stimulus; indexes semantic prediction error &
\href{07_cognitive_integration.md\#sec:cognitive-integration}{§7} \\
P600 & --- & Event-related brain potential peaking \textasciitilde600ms
post-stimulus; indexes syntactic prediction error &
\href{07_cognitive_integration.md\#sec:cognitive-integration}{§7} \\
Perceptual inference & --- & Updating internal beliefs to better predict
current observations (reduce prediction error) &
\href{07_cognitive_integration.md\#sec:cognitive-integration}{§7} \\
Precision weighting & --- & Weighting of prediction errors by inverse
variance; enriched morphism weights serve this role &
\href{07_cognitive_integration.md\#sec:cognitive-integration}{§7} \\
Prediction error & \(\varepsilon\) & Difference between predicted and
observed sensory input; drives belief updating &
\href{07_cognitive_integration.md\#sec:cognitive-integration}{§7} \\
ROSE model & --- & Murphy's
Representation--Operation--Structure--Encoding architecture;
cross-frequency phase-amplitude coupling (PAC) linking biolinguistic
syntax (slow oscillatory phase) to neuropragmatic inference (fast
gamma); the neural-timescale bridge of Pillar 6 &
\href{01_introduction.md\#sec:introduction}{§1},
\href{07c_daif_results.md\#sec:daif-results}{§7c} \\
S-HAI & --- & Schema-Based Hierarchical Active Inference; dual-level
POMDP connecting abstract relational schemas (Level 2: case diagram
structure) to sensorimotor surface parsing (Level 1) &
\href{07_cognitive_integration.md\#sec:cognitive-integration}{§7} \\
Push-forward (general) & \(T_{\#}\mu\) & Image of a measure \(\mu\)
under a measurable map \(T\); DAIF-specific notation in
\autoref{sec:notation-daif} &
\href{07_cognitive_integration.md\#sec:cognitive-integration}{§7} \\
Situation semantics & --- & Framework representing meaning as relations
between situations (Barwise and Perry 1983) &
\href{07_cognitive_integration.md\#sec:cognitive-integration}{§7} \\
Surprisal & \(-\ln p(o)\) & Negative log-probability of an observation
under the generative model &
\href{07_cognitive_integration.md\#sec:cognitive-integration}{§7} \\
Variational free energy & \(F[q]\) & Functional upper bound on surprisal
minimized by approximate posterior \(q\) &
\href{07_cognitive_integration.md\#sec:cognitive-integration}{§7} \\
\bottomrule\noalign{}
\end{longtable}
\end{frame}

\begin{frame}[allowframebreaks]{Quantum and Topological Terms}
\protect\phantomsection\label{sec:notation-quantum}
\begin{longtable}[]{@{}
  >{\raggedright\arraybackslash}p{(\linewidth - 6\tabcolsep) * \real{0.1477}}
  >{\centering\arraybackslash}p{(\linewidth - 6\tabcolsep) * \real{0.0682}}
  >{\raggedright\arraybackslash}p{(\linewidth - 6\tabcolsep) * \real{0.7159}}
  >{\centering\arraybackslash}p{(\linewidth - 6\tabcolsep) * \real{0.0682}}@{}}
\caption{Quantum and topological
terminology.}\label{tbl:not-quantum}\tabularnewline
\toprule\noalign{}
\begin{minipage}[b]{\linewidth}\raggedright
Term
\end{minipage} & \begin{minipage}[b]{\linewidth}\centering
Symbol
\end{minipage} & \begin{minipage}[b]{\linewidth}\raggedright
Definition
\end{minipage} & \begin{minipage}[b]{\linewidth}\centering
First Use
\end{minipage} \\
\midrule\noalign{}
\endfirsthead
\toprule\noalign{}
\begin{minipage}[b]{\linewidth}\raggedright
Term
\end{minipage} & \begin{minipage}[b]{\linewidth}\centering
Symbol
\end{minipage} & \begin{minipage}[b]{\linewidth}\raggedright
Definition
\end{minipage} & \begin{minipage}[b]{\linewidth}\centering
First Use
\end{minipage} \\
\midrule\noalign{}
\endhead
Amplituhedron & --- & Positive geometry encoding scattering amplitudes;
connected to TQNN execution traces &
\href{08_quantum_active_inference.md\#sec:quantum-active-inference}{§8} \\
Barren plateau & --- & Vanishing gradient phenomenon in PQC training;
gradients decay exponentially in system size for certain ansätze &
\href{04b_compact_closure_complexity.md\#sec:compact-closure-complexity}{§4b} \\
\(\dagger\)-compact category & --- & See \(\dagger\)-compact closed
category in \href{11b_notation.md\#sec:notation-category}{Category
Theory} &
\href{08_quantum_active_inference.md\#sec:quantum-active-inference}{§8} \\
CPTP map & --- & Completely Positive Trace-Preserving map; quantum
channel between Hilbert spaces &
\href{08_quantum_active_inference.md\#sec:quantum-active-inference}{§8} \\
Density operator & \(\rho\) & Positive semidefinite trace-one operator
encoding a quantum (or semantic) state &
\href{08_quantum_active_inference.md\#sec:quantum-active-inference}{§8} \\
Execution trace & --- & Record of operations in a quantum computation;
connects TQNNs to amplituhedra &
\href{08_quantum_active_inference.md\#sec:quantum-active-inference}{§8} \\
Generalized flow & --- & Graph-theoretic property ensuring deterministic
circuit extraction from ZX-diagrams &
\href{08_quantum_active_inference.md\#sec:quantum-active-inference}{§8} \\
Hadamard box & \(H\) & ZX-calculus element implementing the Hadamard
gate; converts between Z and X bases &
\href{08_quantum_active_inference.md\#sec:quantum-active-inference}{§8} \\
Holographic screen & --- & Information boundary between interacting
quantum systems carrying a qubit array &
\href{08_quantum_active_inference.md\#sec:quantum-active-inference}{§8} \\
Parameterized quantum circuit (PQC) & --- & Quantum circuit with
trainable angle parameters; the computational substrate of QNLP and
case-category implementations &
\href{04b_compact_closure_complexity.md\#sec:compact-closure-complexity}{§4b} \\
IQP ansatz & --- & Instantaneous Quantum Polynomial-time circuit:
default lambeq PQC ansatz for noun and verb boxes &
\href{04b_compact_closure_complexity.md\#sec:compact-closure-complexity}{§4b} \\
Sim4 ansatz & --- & Strongly entangling layer PQC ansatz used for
discourse-level lambeq Gen II circuits &
\href{04c_discourse_complexity.md\#sec:discocirc-discourse}{§4c} \\
Pointer state & --- & Preferred quantum state selected by a QRF;
determines the measurement basis &
\href{08_quantum_active_inference.md\#sec:quantum-active-inference}{§8} \\
Quantum contextuality & --- & Quantum correlations that reduce
cohomological obstructions to semantic alignment &
\href{08_quantum_active_inference.md\#sec:quantum-active-inference}{§8} \\
Quantum discord & --- & Quantum correlation measure equal to integrated
semantic information in sheaf framework &
\href{08_quantum_active_inference.md\#sec:quantum-active-inference}{§8} \\
Quantum key distribution (QKD) & --- & Protocol providing
information-theoretic security for quantum communication channels &
\href{09_ai_implications.md\#sec:ai-implications}{§9} \\
Quantum reference frame (QRF) & --- & Observer-relative frame selecting
pointer states and inducing decoherence &
\href{08_quantum_active_inference.md\#sec:quantum-active-inference}{§8} \\
Semantic Hilbert space & \(H_v\) & The finite-dimensional semantic
Hilbert space carried at vertex \(v\) in a quantum semantic sheaf &
\href{08b_quantum_semantics.md\#sec:quantum-semantics}{§8b} \\
Quantum semantic sheaf & \((H, F, \rho)\) & Triple of Hilbert spaces,
CPTP channels, and density operators over a communication graph &
\href{08_quantum_active_inference.md\#sec:quantum-active-inference}{§8} \\
Reshetikhin--Turaev invariant & --- & Topological invariant assigning to
a ribbon graph a linear map via TQFT &
\href{08_quantum_active_inference.md\#sec:quantum-active-inference}{§8} \\
Sheaf cohomology & \(H^n(\mathcal{F})\) & Cohomological obstruction
classes governing alignment in a quantum semantic sheaf &
\href{08_quantum_active_inference.md\#sec:quantum-active-inference}{§8} \\
Spin-network & --- & Graph with edges labeled by representations and
vertices by intertwiners; TQFT data &
\href{08_quantum_active_inference.md\#sec:quantum-active-inference}{§8} \\
Spider & --- & Elementary ZX-diagram node (Z-spider or X-spider)
representing a quantum operation &
\href{08_quantum_active_inference.md\#sec:quantum-active-inference}{§8} \\
TQFT & --- & Topological Quantum Field Theory; functor from cobordisms
to Hilbert spaces &
\href{08_quantum_active_inference.md\#sec:quantum-active-inference}{§8} \\
TQNN & --- & Topological Quantum Neural Network; QNN reformulated via
spin-networks and TQFT &
\href{08_quantum_active_inference.md\#sec:quantum-active-inference}{§8} \\
Turaev--Viro invariant & --- & State-sum TQFT invariant; implements
quantum error-correcting codes in TQNNs &
\href{08_quantum_active_inference.md\#sec:quantum-active-inference}{§8} \\
ZX-calculus & --- & Graphical language for quantum circuits using
Z-spiders, X-spiders, and Hadamard boxes &
\href{08_quantum_active_inference.md\#sec:quantum-active-inference}{§8} \\
ZX-diagram & --- & String diagram in a \(\dagger\)-compact closed
category representing a quantum process &
\href{08_quantum_active_inference.md\#sec:quantum-active-inference}{§8} \\
ZX rewrite rule & --- & Graph-theoretic transformation preserving the
semantics (linear map) of a ZX-diagram &
\href{08_quantum_active_inference.md\#sec:quantum-active-inference}{§8} \\
\bottomrule\noalign{}
\end{longtable}
\end{frame}

\begin{frame}[allowframebreaks]{Logical and Type-Theoretic Terms}
\protect\phantomsection\label{sec:notation-logical}
\begin{longtable}[]{@{}
  >{\raggedright\arraybackslash}p{(\linewidth - 6\tabcolsep) * \real{0.1477}}
  >{\centering\arraybackslash}p{(\linewidth - 6\tabcolsep) * \real{0.0682}}
  >{\raggedright\arraybackslash}p{(\linewidth - 6\tabcolsep) * \real{0.7159}}
  >{\centering\arraybackslash}p{(\linewidth - 6\tabcolsep) * \real{0.0682}}@{}}
\caption{Logical and type-theoretic
terminology.}\label{tbl:not-logical}\tabularnewline
\toprule\noalign{}
\begin{minipage}[b]{\linewidth}\raggedright
Term
\end{minipage} & \begin{minipage}[b]{\linewidth}\centering
Symbol
\end{minipage} & \begin{minipage}[b]{\linewidth}\raggedright
Definition
\end{minipage} & \begin{minipage}[b]{\linewidth}\centering
First Use
\end{minipage} \\
\midrule\noalign{}
\endfirsthead
\toprule\noalign{}
\begin{minipage}[b]{\linewidth}\raggedright
Term
\end{minipage} & \begin{minipage}[b]{\linewidth}\centering
Symbol
\end{minipage} & \begin{minipage}[b]{\linewidth}\raggedright
Definition
\end{minipage} & \begin{minipage}[b]{\linewidth}\centering
First Use
\end{minipage} \\
\midrule\noalign{}
\endhead
\(\beta\text{-reduction}\) & --- & Computational reduction step in
\(\lambda\text{-calculus}\); corresponds to cut elimination in proofs &
\href{03_categorial_grammar.md\#sec:categorial-grammar}{§3} \\
Church--Rosser property & --- & Confluence of
\(\beta\text{-reduction}\): all reduction sequences converge to the same
normal form &
\href{03_categorial_grammar.md\#sec:categorial-grammar}{§3} \\
Curry--Howard isomorphism & --- & Correspondence between proofs and
programs, propositions and types &
\href{03_categorial_grammar.md\#sec:categorial-grammar}{§3} \\
Cut elimination & --- & Proof normalization procedure removing
intermediate lemmas; corresponds to \(\beta\text{-reduction}\) &
\href{03_categorial_grammar.md\#sec:categorial-grammar}{§3} \\
Graded type theory & --- & Extension of type theory tracking effects
(e.g., evidentiality) via graded modalities &
\href{03_categorial_grammar.md\#sec:categorial-grammar}{§3} \\
Lambek calculus & --- & Non-commutative intuitionistic linear logic for
syntactic type assignment &
\href{03_categorial_grammar.md\#sec:categorial-grammar}{§3} \\
Left residual & \(A \backslash B\) & Type of an expression that, given
\(A\) to the left, produces \(B\) &
\href{03_categorial_grammar.md\#sec:categorial-grammar}{§3} \\
Residuation law &
\(A \otimes B \leq C \iff A \leq C / B \iff B \leq A \backslash C\) &
Fundamental axiom connecting the three connectives of the Lambek
calculus &
\href{03_categorial_grammar.md\#sec:categorial-grammar}{§3} \\
Right residual & \(B / A\) & Type of an expression that, given \(A\) to
the right, produces \(B\) &
\href{03_categorial_grammar.md\#sec:categorial-grammar}{§3} \\
\bottomrule\noalign{}
\end{longtable}
\end{frame}

\begin{frame}[allowframebreaks]{AI and Communication Protocols}
\protect\phantomsection\label{sec:notation-ai}
\begin{longtable}[]{@{}
  >{\raggedright\arraybackslash}p{(\linewidth - 6\tabcolsep) * \real{0.1477}}
  >{\centering\arraybackslash}p{(\linewidth - 6\tabcolsep) * \real{0.0682}}
  >{\raggedright\arraybackslash}p{(\linewidth - 6\tabcolsep) * \real{0.7159}}
  >{\centering\arraybackslash}p{(\linewidth - 6\tabcolsep) * \real{0.0682}}@{}}
\caption{AI and communication protocol
terminology.}\label{tbl:not-ai-protocols}\tabularnewline
\toprule\noalign{}
\begin{minipage}[b]{\linewidth}\raggedright
Term
\end{minipage} & \begin{minipage}[b]{\linewidth}\centering
Symbol
\end{minipage} & \begin{minipage}[b]{\linewidth}\raggedright
Definition
\end{minipage} & \begin{minipage}[b]{\linewidth}\centering
First Use
\end{minipage} \\
\midrule\noalign{}
\endfirsthead
\toprule\noalign{}
\begin{minipage}[b]{\linewidth}\raggedright
Term
\end{minipage} & \begin{minipage}[b]{\linewidth}\centering
Symbol
\end{minipage} & \begin{minipage}[b]{\linewidth}\raggedright
Definition
\end{minipage} & \begin{minipage}[b]{\linewidth}\centering
First Use
\end{minipage} \\
\midrule\noalign{}
\endhead
A2A Protocol & --- & Google's Agent-to-Agent protocol for
cross-framework agent communication via HTTP/JSON-RPC &
\href{09_ai_implications.md\#sec:ai-implications}{§9} \\
ACP & --- & Agent Communication Protocol; standardizes messaging formats
across agents, apps, and humans &
\href{09_ai_implications.md\#sec:ai-implications}{§9} \\
ANP & --- & Agent Network Protocol; three-layer architecture for trusted
distributed agent interaction &
\href{09_ai_implications.md\#sec:ai-implications}{§9} \\
Categorical deep learning & --- & Deep learning approached through the
lens of category theory (Gavranović et al.) &
\href{09_ai_implications.md\#sec:ai-implications}{§9} \\
Double Categorical Systems Theory (DCST) & --- & Framework using
2-categories (horizontal + vertical composition) for explainable
autonomous AI & \href{09_ai_implications.md\#sec:ai-implications}{§9} \\
Functorial encryption & --- & Semantic cryptography: applying a secret
functor to map plaintext categories into ciphertext categories &
\href{09_ai_implications.md\#sec:ai-implications}{§9} \\
lambeq & --- & Quantum Natural Language Processing pipeline compiling
DisCoCat diagrams to quantum circuits &
\href{09_ai_implications.md\#sec:ai-implications}{§9} \\
MCP & --- & Model Context Protocol; standardizes how AI agents access
external tools and data sources &
\href{09_ai_implications.md\#sec:ai-implications}{§9} \\
Parameterized optics / lenses & --- & Categorical constructions modeling
neural network components (Gavranović); attention heads as optics &
\href{09_ai_implications.md\#sec:ai-implications}{§9} \\
QNLP & --- & Quantum Natural Language Processing; quantum computation on
DisCoCat sentence diagrams &
\href{09_ai_implications.md\#sec:ai-implications}{§9} \\
Role variables & \(X_{\text{NOM}}, X_{\text{ACC}}, \ldots\) & Agentic
components structured by case roles in networked LLM contexts (e.g.,
active requester policy) &
\href{09_ai_implications.md\#sec:ai-implications}{§9} \\
Semantic cryptography & --- & Encrypting compositional meaning
structures (functorial encryption, diagram obfuscation, weight masking)
& \href{09_ai_implications.md\#sec:ai-implications}{§9} \\
Protocol category & \(\mathcal{C}_{\text{protocol}}\) & Category whose
morphisms are licensed interaction steps in the case-theoretic firewall
(\autoref{sec:cognitive-security}) &
\href{09b_cognitive_security.md\#sec:cognitive-security}{§9b} \\
Adversarial morphism & \(\phi{}\) & Illicit map attempting case-role
promotion (e.g.~ACC\(\to\)NOM) in injection attacks &
\href{09b_cognitive_security.md\#sec:cognitive-security}{§9b} \\
Access Collapse & --- & Catastrophic boundary failure where adversarial
text traverses from passive data (ACC) to active instruction (NOM)
without structural constraint &
\href{09b_cognitive_security.md\#sec:cognitive-security}{§9b} \\
Case-theoretic firewall & --- & Type-checking system enforcing licensed
morphism constraints at agent communication boundaries; detects illicit
role promotions in \(\text{Mor}(\mathcal{C}_{\text{protocol}})\) &
\href{09b_cognitive_security.md\#sec:cognitive-security}{§9b} \\
Prompt injection & --- & Attack illicitly promoting user-supplied data
from ACC (patient) to NOM (commanding agent); structurally a type
violation in the interaction grammar &
\href{09b_cognitive_security.md\#sec:cognitive-security}{§9b} \\
Symbolic Isolation & --- & Property ensuring ACC-typed data wires cannot
fuse with NOM-typed command wires; enforced by the non-cartesian
fragment of the monoidal structure &
\href{09b_cognitive_security.md\#sec:cognitive-security}{§9b} \\
\bottomrule\noalign{}
\end{longtable}
\end{frame}

\begin{frame}[fragile,allowframebreaks]{Diagrammatic Reasoning}
\protect\phantomsection\label{sec:notation-diagrammatic}
\begin{longtable}[]{@{}
  >{\raggedright\arraybackslash}p{(\linewidth - 6\tabcolsep) * \real{0.1477}}
  >{\centering\arraybackslash}p{(\linewidth - 6\tabcolsep) * \real{0.0682}}
  >{\raggedright\arraybackslash}p{(\linewidth - 6\tabcolsep) * \real{0.7159}}
  >{\centering\arraybackslash}p{(\linewidth - 6\tabcolsep) * \real{0.0682}}@{}}
\caption{Diagrammatic reasoning
terminology.}\label{tbl:not-diagrammatic}\tabularnewline
\toprule\noalign{}
\begin{minipage}[b]{\linewidth}\raggedright
Term
\end{minipage} & \begin{minipage}[b]{\linewidth}\centering
Symbol
\end{minipage} & \begin{minipage}[b]{\linewidth}\raggedright
Definition
\end{minipage} & \begin{minipage}[b]{\linewidth}\centering
First Use
\end{minipage} \\
\midrule\noalign{}
\endfirsthead
\toprule\noalign{}
\begin{minipage}[b]{\linewidth}\raggedright
Term
\end{minipage} & \begin{minipage}[b]{\linewidth}\centering
Symbol
\end{minipage} & \begin{minipage}[b]{\linewidth}\raggedright
Definition
\end{minipage} & \begin{minipage}[b]{\linewidth}\centering
First Use
\end{minipage} \\
\midrule\noalign{}
\endhead
Categorical complexity & \(\kappa(D)\) & Complexity metric for diagram
\(D\) derived from box counts, cup counts, and swap depths &
\href{04b_compact_closure_complexity.md\#sec:compact-closure-complexity}{§4b} \\
Cognitively privileged representation & --- & Representation format that
leverages perceptual and spatial cognition for inference &
\href{01_introduction.md\#sec:introduction}{§1} \\
Diagram depth & --- & Length of the longest input-to-output path through
boxes; measures derivational complexity &
\href{04_categorical_semantics.md\#sec:categorical-semantics}{§4} \\
Existential graphs & --- & Peirce's graphical logic system conducting
first-order logic entirely diagrammatically &
\href{07_cognitive_integration.md\#sec:cognitive-integration}{§7} \\
Free ride & --- & Shimojima's term: information extracted from a diagram
without explicit inference steps &
\href{01_introduction.md\#sec:introduction}{§1} \\
Hybrid reasoning & --- & Giardino's term: reasoning combining perceptual
pattern recognition with theoretical knowledge &
\href{01_introduction.md\#sec:introduction}{§1} \\
Inferential instrument & --- & Manders's term: a diagram whose spatial
properties encode proof-relevant information &
\href{06_topos_theory.md\#sec:topos-theory}{§6} \\
Joyal--Street theorem & --- & Soundness and completeness of
string-diagrammatic reasoning for monoidal categories &
\href{03_categorial_grammar.md\#sec:categorial-grammar}{§3} \\
Normal form & \(D_{\text{nf}}\) & Canonical form of a diagram obtained
by rewriting; unique up to the axioms &
\href{04_categorical_semantics.md\#sec:categorical-semantics}{§4} \\
Role colours (figures) & \texttt{CASE\_COLORS} & \texttt{CaseRole} →
display colour for generated figures; single source in
\texttt{src/visualization/styles.py} &
\href{02_case_systems.md\#sec:case-systems}{§2}--\href{04_categorical_semantics.md\#sec:categorical-semantics}{§4},
\href{11_syntactic_sentence_diagrams.md\#sec:syntactic-diagrams}{Appendix
A} \\
\bottomrule\noalign{}
\end{longtable}
\end{frame}

\begin{frame}[fragile,allowframebreaks]{Notation Conventions}
\protect\phantomsection\label{sec:notation-conventions}
\begin{longtable}[]{@{}
  >{\raggedright\arraybackslash}p{(\linewidth - 2\tabcolsep) * \real{0.1286}}
  >{\raggedright\arraybackslash}p{(\linewidth - 2\tabcolsep) * \real{0.8714}}@{}}
\caption{General mathematical notation and manuscript
conventions.}\label{tbl:not-conventions}\tabularnewline
\toprule\noalign{}
\begin{minipage}[b]{\linewidth}\raggedright
Convention
\end{minipage} & \begin{minipage}[b]{\linewidth}\raggedright
Meaning
\end{minipage} \\
\midrule\noalign{}
\endfirsthead
\toprule\noalign{}
\begin{minipage}[b]{\linewidth}\raggedright
Convention
\end{minipage} & \begin{minipage}[b]{\linewidth}\raggedright
Meaning
\end{minipage} \\
\midrule\noalign{}
\endhead
\(\mathcal{C}, \mathcal{D}\) & Categories \\
\(\mathcal{V}\) & Enrichment base (monoidal category, typically
\(([0,1], \cdot, 1)\)) \\
\(\mathcal{U}\) & Universal (maximal) case category \\
\(\mathcal{L}\) & Language-specific case category \\
\(\mathcal{E}_{\mathbb{T}}\) & Classifying topos of theory
\(\mathbb{T}\) \\
\(f, g, h\) & Morphisms \\
\(F, G\) & Functors \\
\(\alpha{}\), \(\beta{}\) & Natural transformations; functorial vertical
composition (componentwise along objects) is standard
(\href{02_case_systems.md\#sec:case-systems}{§2}). In DAIF,
\(\alpha_{\mathrm{pol}}\) is policy temperature;
\(\alpha_{\mathrm{CVaR}}\) is CVaR scale; \(\beta_{\mathrm{risk}}\) is
EFE risk weight; see
\href{07c_daif_results.md\#sec:daif-results}{§7c} \\
\(n, s\) & Basic pregroup types: noun, sentence \\
\(n^l, n^r\) & Left and right adjoints of type \(n\) \\
\(n_{\text{NOM}}, n_{\text{ACC}}\) & Case-subscripted noun types (e.g.,
nominative noun, accusative noun) \\
\(N, S\) & Noun space and sentence space under the meaning functor \\
\(N_{\text{NOM}}, N_{\text{ACC}}, \ldots\) & Case-specific vector
subspaces in case-enriched DisCoCat \\
\(\overrightarrow{\text{word}}\) & Word vector (column vector in noun
space \(N\)) \\
\(\overleftrightarrow{\text{verb}}\) & Verb tensor (element of
\(N \otimes S \otimes N\) for transitive verbs) \\
\(\mathbf{Preg}\) & Category of pregroup types and reductions \\
\(\mathbf{FVect}\) / \(\mathbf{FdVect}\) & Category of
finite-dimensional vector spaces and linear maps \\
\(\mathbf{FHilb}\) & Category of finite-dimensional Hilbert spaces and
linear maps \\
\(\mathbf{Qubit}\) & Category of qubit systems with tensor product
structure \\
\(\mathbf{Set}\) & Category of sets and functions \\
\(\mathbf{Ent}, \mathbf{Case}\) & Categories for Entities and Case Roles
respectively in DisCoCirc discourse extensions \\
\(\otimes\) & Tensor product (monoidal product, type concatenation) \\
\(\circ\) & Composition of morphisms \\
\(\Rightarrow\) & Natural transformation between functors \\
\(\simeq\) & Categorical equivalence \\
\(\leq\) & Preorder relation on types (derivability) \\
\(\gamma{}\) & Discount factor in distributional RL return
computation \\
\(w \in [0,1]\) & Enriched morphism weight (proto-role satisfaction
degree) \\
\(\varepsilon\) & Sensory prediction error (active inference);
\(\varepsilon_n\) also denotes the compact-closure \textbf{counit} (cup)
on type \(n\) \\
\(\eta_{\mathrm{cc}}\) & Compact closure \textbf{unit} (cap); IQN
curvature uses \(\eta_{\mathrm{IQN}}\) instead \\
\(\epsilon\) & Small numerical floor (e.g.~KL stabiliser, VMP
convergence threshold); not the counit \\
\(\rho{}\) & Density operator (quantum state) \\
\([@key]\) & Parenthetical citation \\
\texttt{{[}-@key{]}} & Suppress-author citation \\
\texttt{\textbackslash{}autoref} (with \texttt{\{sec:…\}} /
\texttt{\{eq:…\}} / \texttt{\{fig:…\}}) & LaTeX/Markdown automatic
cross-reference to a labeled target \\
\(Z(s,a)\) & Return distribution (DAIF); full distributional
representation of returns \\
\(G(\pi)\) & Expected free energy of policy (DAIF active inference) \\
\(\tau{}\) & Quantile level in QR-DQN / IQN \\
\(\psi_{\mathrm{IQN}}(\tau)\) & IQN risk distortion of quantile
levels \\
\(\beta_{\mathrm{risk}}\) & Risk-sensitivity coefficient in
\(G(\pi)\) \\
\(\alpha_{\mathrm{pol}} \;;\; \alpha_{\mathrm{CVaR}}\) & Policy softmax
temperature; CVaR tail level (IQN) \\
\(w_f, w_k, w_c\) & Enriched morphism / role weights (precision on PE
and VMP) \\
\(\mathrm{DPE}\) & Distributional prediction error; precision-weighted
Wasserstein mismatch \\
\(H[Z]\) & Return distribution entropy \\
\bottomrule\noalign{}
\end{longtable}
\end{frame}

\end{document}
